\documentclass[conference]{IEEEtran}

\usepackage[T1]{fontenc}
% Courier w T1 nie ma w BasicTeX; Latin Modern jest w każdej dystrybucji.
\renewcommand{\ttdefault}{lmtt}
\usepackage{amsmath,amssymb}
\usepackage{graphicx}
\usepackage{booktabs}
\usepackage{siunitx}
\usepackage{listings}
\usepackage[hidelinks]{hyperref}
\usepackage[capitalise,nameinlink]{cleveref}

\graphicspath{{../results/figures/}}

\begin{document}

\title{Privacy-Preserving Spam Detection\\on Data Encrypted with the Paillier Cryptosystem}

\author{
  \IEEEauthorblockN{Krzysztof Śliwiński, Piotr Waluszek}
  \IEEEauthorblockA{Faculty of Computer Science and Telecommunications\\
    Cracow University of Technology, Kraków, Poland}
}

\maketitle

% Każda sekcja w osobnym pliku, żeby dwie osoby mogły pisać równolegle bez konfliktów.
% P23
\begin{abstract}
TODO
\end{abstract}

\begin{IEEEkeywords}
homomorphic encryption, Paillier cryptosystem, privacy-preserving machine learning, naive Bayes, spam detection
\end{IEEEkeywords}

% P23
\section{Introduction}\label{sec:introduction}

Spam filtering is commonly offered as a cloud service: the provider maintains a model trained on large amounts of labelled data, and clients send their messages to be classified.
The price is privacy, since the provider sees the content of every message, including the legitimate ones, which are the large majority.
Homomorphic encryption removes this trade-off for models that can be computed on encrypted data: the client sends its message encrypted, the server computes the classifier's score without decrypting it, and only the client learns the result.

Fully homomorphic schemes can evaluate arbitrary functions but are expensive.
Many classifiers, however, only need a linear function of the input, which an additively homomorphic scheme such as Paillier~\cite{paillier1999} computes directly.
Multinomial naive Bayes is one of them: for two classes, its decision is the sign of a log-odds score that is an inner product of the features with fixed weights plus a constant (\cref{sec:background}).
Bost et al.~\cite{bost2015} protect both the model and the data and need interactive comparison protocols for that; we keep the model public to the server, as in a typical cloud service, which reduces classification to a single round trip with one encrypted score.

In this work, we build and evaluate such a privacy-preserving spam filter for SMS messages.
Our contributions are:
\begin{itemize}
  \item a protocol in which the server computes the naive Bayes log-odds on Paillier-encrypted TF-IDF features and returns one ciphertext per message, with a fixed-point encoding that makes the encrypted computation exact and hides which features are non-zero (\cref{sec:system-design});
  \item an evaluation on the SMS Spam Collection~\cite{smsspam} with duplicates removed, showing that the encrypted classifier reaches the same quality as its plaintext counterpart, \SI{98.07}{\percent} accuracy, with identical predictions and bit-exact scores (\cref{sec:setup,sec:results});
  \item a cost analysis for key lengths from \num{1024} to \num{3072} bits and $500$ to $1000$ features, which together with the accuracy for $100$ to $1000$ features motivates the number of features (\cref{sec:results});
  \item an empirical security analysis against an honest-but-curious server that can encrypt chosen messages, which shows that a classifier on ciphertexts does not beat the majority-class baseline, while encrypting only the non-zero features would leak almost as much as the plaintext, and reusing the encryption randomness or the floating-point encoding of \texttt{phe} would let the server read features directly (\cref{sec:threat-model,sec:security}).
\end{itemize}

The implementation uses the \texttt{phe} library~\cite{phe} and scikit-learn, and all results can be reproduced with a single command.

% K18
\section{Background}\label{sec:background}

TODO

% K19
\section{Threat Model}\label{sec:threat-model}

\textit{Setting.}
A client wants its messages classified by a spam filter hosted by a cloud provider without revealing their content.
The provider trains the naive Bayes model on its own plaintext data (in our experiments, the training part of the SMS Spam Collection~\cite{smsspam}); encryption protects only inference, i.e.\ the client's messages.

\textit{Adversary profile.}
The cloud server is \emph{honest-but-curious}: a passive adversary that follows the protocol and returns correct results, but records and analyses everything it observes in order to learn about the messages or their labels.

\textit{Adversary interface.}
The server observes:
\begin{itemize}
  \item the client's public key $n$;
  \item for every message, $k$ ciphertexts $\mathrm{Enc}(x_i)$, one per feature, including features equal to zero, together with the public exponent attached by \texttt{phe} (constant for all values in our design, see \cref{sec:system-design});
  \item the model $(b, \mathbf{w})$, the vocabulary and all TF-IDF parameters, since it trained the model;
  \item metadata: the number of messages, their arrival times and the time of each request.
\end{itemize}
It never sees the private key or any decrypted result.
Since the server holds the public key, it can encrypt any plaintext of its choice, so the relevant attack model is chosen-plaintext (CPA) rather than ciphertext-only.
The security goal is that the ciphertexts reveal nothing about the messages beyond the metadata above, which is exactly the IND-CPA guarantee of Paillier under DCRA.

\textit{Kerckhoffs' principle.}
The protocol, source code, model, vocabulary and feature extraction are public.
The only secret is the client's private key.

\textit{Out of scope.}
An active (malicious) server that returns wrong scores, timing side channels on the client, and confidentiality of the model against the client.
The client learns $\Delta$ for each of its messages, which leaks information about $(b, \mathbf{w})$ and could enable model extraction; protecting the model is not a goal of this work.

\textit{Success criteria.}
The encrypted classifier is successful if:
\begin{enumerate}
  \item its accuracy on the test set exceeds \SI{94}{\percent};
  \item it predicts the same label as the plaintext classifier for \SI{100}{\percent} of test messages;
  \item the decrypted score is exactly equal to the score computed on the same fixed-point integers in plaintext.
\end{enumerate}
The attacks in \cref{sec:security} succeed if a classifier trained on information available to the server predicts spam better than the majority-class baseline.

% K20
\section{System Design}\label{sec:system-design}

\subsection{Protocol}

The client holds a Paillier key pair; the server holds the public key and the model $(b, \mathbf{w})$ from \cref{eq:log-odds}.
Classifying one message takes a single round trip (\cref{tab:protocol}).
Instead of scoring the two classes separately, the server returns one encrypted log-odds value: this halves the server's work and the number of decryptions, gives the score used for ROC analysis directly, and reveals less of the model to the client.

\begin{table}[htbp]
  \centering
  \caption{One protocol run for a single message. $S$ is the fixed-point scale (\cref{sec:fixed-point}).}
  \label{tab:protocol}
  \begin{tabular}{@{}l l p{0.62\columnwidth}@{}}
    \toprule
    Step & Party & Operation \\
    \midrule
    1 & Client & Extract TF-IDF features $\mathbf{x} \in [0, 1]^k$ with the public vocabulary; encode $\tilde{x}_i = \lfloor S x_i \rceil$; send $c_i = \mathrm{Enc}(\tilde{x}_i)$ for \emph{all} $i = 1, \dots, k$. \\
    2 & Server & Encode $\tilde{w}_i = \lfloor S w_i \rceil$, $\tilde{b} = \lfloor S^2 b \rceil$; compute $c_\Delta = \mathrm{Enc}(\tilde{b}) \cdot \prod_{i} c_i^{\tilde{w}_i} \bmod n^2$ and return it. \\
    3 & Client & Decrypt $\tilde{\Delta} = \mathrm{Dec}(c_\Delta)$; label the message as spam if and only if $\tilde{\Delta} > 0$; use $\tilde{\Delta} / S^2$ as the score. \\
    \bottomrule
  \end{tabular}
\end{table}

By \cref{eq:hom-add,eq:hom-mul}, $c_\Delta$ is an encryption of $\tilde{\Delta} = \tilde{b} + \sum_i \tilde{w}_i \tilde{x}_i$, the log-odds computed on the encoded integers.
Negative weights need no special handling: \texttt{phe} represents them modulo $n$ and evaluates $c_i^{\tilde{w}_i}$ through a modular inverse.

\subsection{Fixed-Point Encoding}\label{sec:fixed-point}

Features and weights are scaled by $S = 10^{6}$ and rounded to integers; the bias is scaled by $S^2$, so that it has the same scale as the products $\tilde{w}_i \tilde{x}_i$.
We encode values ourselves instead of relying on the floating-point encoding of \texttt{phe} for two reasons.

\textit{No leakage through the exponent.}
For a float, \texttt{phe} chooses the public exponent attached to the ciphertext from the magnitude of the value, so it is not the same for every value: zero shares its exponent with values in $[1/16, 1)$, but smaller values and $1.0$ get different ones.
The server, which knows the encoding, would learn that the corresponding words occur in the message (\cref{sec:security}).
An integer always gets exponent~0, so all ciphertexts look alike.

\textit{Exact, verifiable computation.}
Over the integers, the homomorphic computation is exact: the decrypted $\tilde{\Delta}$ must be \emph{equal} to the same integer computation done in plaintext, which is our third success criterion (\cref{sec:threat-model}).
The only approximation is the rounding itself.
A zero feature is encoded exactly and contributes nothing, so the error of the score is bounded by a sum over the non-zero features only:
\begin{equation}
  \left| \frac{\tilde{\Delta}}{S^2} - \Delta \right| \le \frac{1}{2S} \sum_{i:\, x_i \ne 0} \left( |x_i| + |w_i| + \frac{1}{2S} \right) + \frac{1}{2S^2}.
\end{equation}
For SMS messages with a dozen non-zero features this is of the order of $10^{-5}$; we measure the actual error against the floating-point classifier in \cref{sec:results}.
The encoded values are also far from the plaintext limit: $|\tilde{\Delta}| \le S^2 (|b| + \sum_i |w_i|) < 10^{17}$, whereas \texttt{phe} accepts signed values up to $n/3 \approx 2^{2046}$ for a \num{2048}-bit key.

\subsection{Full Feature Vectors}

The client encrypts all $k$ features, including the zeros.
Sending only the non-zero positions would be cheaper but would reveal which vocabulary words occur in each message; we show in \cref{sec:security} that this alone is enough to classify spam.
The price is bandwidth: with \SI{512}{\byte} per ciphertext, a message costs \SI{256}{\kilo\byte} for $k = 500$ and \SI{512}{\kilo\byte} for $k = 1000$.

\subsection{Implementation}

The system is written in Python, and the model is trained with scikit-learn.
The weights $b$ and $\mathbf{w}$ are the differences between the spam and ham rows of \texttt{class\_log\_prior\_} and \texttt{feature\_log\_prob\_} of the fitted \texttt{MultinomialNB}, respectively.
The cryptographic operations use \texttt{phe} with \texttt{gmpy2}.
Client-side encryption ($k$ full encryptions per message) dominates the cost, so it runs in parallel across processes.
The server performs $k$ exponentiations with small exponents per message, and the client a single decryption.

% P20
\section{Experimental Setup}\label{sec:setup}

\textit{Data.}
We use the SMS Spam Collection~\cite{smsspam}: \num{5574} English SMS messages, \num{747} of them spam.
The collection contains \num{403} exact duplicates of message texts (\num{309} ham, \num{94} spam), typically mass-sent notifications and spam campaigns.
We keep the first occurrence of each text, which leaves \num{5171} messages with \num{653} spam (\SI{12.6}{\percent}).
Without deduplication, \SI{12}{\percent} of the test messages (and \SI{22}{\percent} of the spam among them) would have an identical copy in the training set, which inflates the measured accuracy.

\textit{Split.}
The deduplicated data are split once into a training set of \num{4136} messages (\num{522} spam) and a test set of \num{1035} messages (\num{131} spam), stratified by class with a fixed seed.
All model selection uses the training set only; the test set is used once, for the final results in \cref{sec:results}.

\textit{Model and hyperparameters.}
The plaintext pipeline consists of a TF-IDF vectorizer (lowercased tokens of at least two word characters, L2 normalization), selection of the $k$ features with the highest $\chi^2$ score, and multinomial naive Bayes with additive smoothing~$\alpha$.
TF-IDF and the $\chi^2$ selection are part of the pipeline, so in cross-validation they are refitted on every training fold and do not see the validation fold.
We search a grid of 168 configurations: unigrams or unigrams and bigrams, raw or sublinear term frequency, $k \in \{100, 200, 300, 500, 750, 1000\}$ and $\alpha \in \{0.001, 0.003, 0.01, 0.03, 0.1, 0.3, 1\}$, with stratified 5-fold cross-validation on the training set.
We cap $k$ at \num{1000} because the client's cost grows linearly with $k$.
Since \SI{87}{\percent} of the messages are ham, accuracy separates the models poorly, so we select the configuration with the highest F1 score of the spam class.
The selected model uses unigrams, sublinear term frequency, $k = 1000$ and $\alpha = 0.01$, with a cross-validated accuracy of $0.980 \pm 0.008$ and F1 of $0.915 \pm 0.036$.
Its weights $(b, \mathbf{w})$ are exported as described in \cref{sec:system-design}.

\textit{Baselines.}
We compare the encrypted classifier with three control models evaluated on the same test set:
\begin{itemize}
  \item the \emph{plaintext} naive Bayes model, i.e.\ the same model evaluated in double-precision floating point without encryption or fixed-point encoding; this is the reference for the agreement criteria;
  \item a \emph{majority-class} classifier that always predicts ham, which shows how much accuracy the class imbalance alone gives;
  \item \emph{logistic regression} on the same features, with the regularization strength and class weighting tuned by the same cross-validation ($C = 100$, balanced class weights); it is also a linear model $b + \mathbf{x} \cdot \mathbf{w}$, so it could be evaluated with the same protocol.
\end{itemize}

\textit{Metrics.}
We report accuracy and the precision, recall and F1 score of the spam class.
ROC-AUC is computed from the continuous score: the decrypted log-odds $\tilde{\Delta} / S^2$ for the encrypted classifier and $\Delta$ for the plaintext one.
Agreement is measured as the fraction of test messages with the same label as the plaintext model and the fraction whose decrypted integer equals the plaintext fixed-point computation (\cref{sec:threat-model}).

\textit{Implementation and hardware.}
All experiments run on an Apple M5 Pro (15 cores, \SI{24}{\giga\byte} RAM) with Python~3.14, scikit-learn~1.9 and \texttt{phe}~1.5~\cite{phe} accelerated by \texttt{gmpy2}.
The final results use a \num{2048}-bit modulus $n$ and the scale $S = 10^6$; client-side encryption of the test set is parallelized over 15 processes.
The cost per message is measured separately on a single core, for \num{1024}-, \num{2048}- and \num{3072}-bit keys and $k \in \{500, 750, 1000\}$, averaged over 10 messages.
Random seeds are fixed, so all metrics are reproducible; ciphertexts are not, since \texttt{phe} draws a fresh random $r$ for every encryption.

% P21
\section{Results}\label{sec:results}

TODO

% P22 (ataki 1--2), K22 (ataki 3--4)
\section{Security Analysis}\label{sec:security}

We test the threat model of \cref{sec:threat-model} empirically.
Each attack trains a classifier on the information available to the honest-but-curious server and succeeds if it predicts spam better than the majority-class baseline.
Since the server holds the public key and knows the feature extraction, it can prepare labelled examples itself by encrypting messages of its choice (CPA); we let it use the training set for this and never the test set.

\subsection{Attack 1: Classifier on Ciphertexts}\label{sec:attack-ciphertext}

The server encrypts \num{400} labelled training messages (\num{50} spam, stratified sample) with the client's \num{2048}-bit public key, exactly as the client would, and tries to learn spam from the ciphertexts.
From each of the $k = 1000$ ciphertexts $c$ of a message it extracts six features that could reveal structure of the plaintext: the least significant byte, $c \bmod 3$, $c \bmod 5$, $c \bmod 7$, the number of leading zero bits relative to $n^2$, and the Jacobi symbol $\left(\frac{c}{n}\right)$.
The last one breaks the Goldwasser--Micali cryptosystem~\cite{goldwasser1984}, but here $\left(\frac{c}{n}\right) = \left(\frac{r}{n}\right)$, because $(1 + n)^m \equiv 1 \pmod{n}$; it depends only on the random $r$.
A random forest with balanced class weights is trained on the resulting \num{6000} features and evaluated with stratified 5-fold cross-validation.

As a positive control, the same random forest is trained on the plaintext features of the same \num{400} messages.
It reaches an accuracy of $0.975 \pm 0.016$ and an F1 score of $0.887 \pm 0.080$, so the sample is large enough to learn spam if the ciphertexts carried any information about it.
On the ciphertexts, the forest predicts ham for every message: its accuracy of $0.875$ and F1 of $0$ are exactly those of the majority-class baseline (\cref{tab:attacks}).
The attack fails, as IND-CPA predicts.

\subsection{Attack 2: Sparse Feature Vectors}\label{sec:attack-sparse}

A tempting optimization is to encrypt and send only the non-zero features together with their indices.
After the $\chi^2$ selection a message has a median of only 3 non-zero features out of \num{1000}, so this would cut the client's cost and the bandwidth by more than two orders of magnitude.
The values would remain encrypted, but the server would learn \emph{which} vocabulary words occur in every message.

To measure what this reveals, we give the server the binary vector of non-zero positions.
No decryption is needed: the server sees these positions in plaintext.
A Bernoulli naive Bayes model with default parameters, trained on the presence vectors of the training set, classifies the test set with an accuracy of \SI{98.4}{\percent} and F1 of $0.931$, as well as the plaintext model itself (\cref{tab:attacks}).
Even if the indices were hidden and only the number of ciphertexts per message were visible, a logistic regression on this single number reaches \SI{94.3}{\percent}, well above the \SI{87.3}{\percent} of the majority class, because the $\chi^2$ selection favours spam-indicative words, so spam messages contain more of them.
Hence the client must encrypt and send all $k$ features, including the zeros (\cref{sec:system-design}); the protocol then reveals nothing beyond the fixed length $k$.

\begin{table}[htbp]
  \centering
  \caption{Attacks 1 and 2. Attack~1: stratified 5-fold cross-validation on 400 training messages; attack~2: trained on the training set and evaluated on the test set. F1 refers to the spam class.}
  \label{tab:attacks}
  \begin{tabular}{@{}llcc@{}}
    \toprule
    & Information used by the server & Acc. & F1 \\
    \midrule
    1 & Ciphertexts (random forest) & 0.875 & 0.000 \\
      & Majority class & 0.875 & 0.000 \\
      & Plaintext features (control) & 0.975 & 0.887 \\
    \midrule
    2 & Non-zero positions (Bernoulli NB) & 0.984 & 0.931 \\
      & Number of non-zero positions & 0.943 & 0.745 \\
      & Plaintext NB (reference) & 0.981 & 0.922 \\
      & Majority class & 0.873 & 0.000 \\
    \bottomrule
  \end{tabular}
\end{table}

% K22: podsekcje dla ataków 3 (niedeterminizm) i 4 (wykładnik).

% K23
\section{Conclusions}\label{sec:conclusions}

TODO


% Literatura z karty projektu: w bibliografii zawsze, nawet zanim pojawią się \cite w tekście.
\nocite{paillier1999,phe,smsspam}

\bibliographystyle{IEEEtran}
\bibliography{references}

\end{document}
